\section{Introduction and setup}\label{sec:setup}

The static (leakage) current of a CMOS circuit depends strongly on process parameters, supply voltage and temperature (PVT). Estimating it with
full SPICE simulation for every corner and every input vector is expensive. This project replaces SPICE with machine learning in a hierarchical way:
\begin{enumerate}
  \item \textbf{Part A -- transistor models.} Single NMOS and PMOS transistors are simulated in their ON and OFF states over random PVT and bias
        conditions, and one ML model per case predicts the four terminal currents $(I_d, I_g, I_s, I_b)$.
  \item \textbf{Part B -- gates.} Using only these four models, the leakage of INV, NAND2, NOR2, AND2 and OR2 is predicted for every input combination
        (18 cases) at 200 random PVT corners and compared with transistor-level SPICE simulations of the gates.
  \item \textbf{Part C -- test circuit.} The leakage of the assigned test circuit is composed from the gate-level predictions and compared with a SPICE
        simulation of the complete netlist at the same corners.
\end{enumerate}

\subsection{Assignment and tools}
\begin{center}\small
\begin{tabular}{>{\bfseries}l p{0.62\textwidth}}
\toprule
Item & Value\\
\midrule
Team & 26 -- Return of LUZAAZ\\
Model file & PTM \texttt{22nm\_LP.pm} (low-power 22\,nm, BSIM4 \texttt{level=54}, \texttt{version=4.0}), NMOS and PMOS cards\\
Nominal supply & $\Vdd=1\,\mathrm{V}$, sampled uniformly in $0.9$--$1.1\,\mathrm{V}$\\
Channel length & $L=\Lmin=22\,\mathrm{nm}$; widths $W=1$--$8\times\Lmin$ for NMOS and PMOS\\
Test circuit & 8-bit parity checker\\
Simulator & ngspice-47 (official 64-bit Windows build, KLU solver), DC operating point (\texttt{op}) analyses in batch mode\\
Solver options & \texttt{.options gmin=1e-18 abstol=1e-18 reltol=1e-6 vntol=1e-9} in every netlist\\
Software & Python 3.12 with NumPy, pandas, scikit-learn, LightGBM and PyTorch; exact versions locked in \texttt{uv.lock}\\
Randomness & every random draw uses fixed seeds from \texttt{config/project.yaml} (base seed 20261105)\\
\bottomrule
\end{tabular}
\end{center}

Following a clarification in the course group, the drain and source junction areas are set to $AD=AS=W\cdot L_{\mathrm{diff}}$ with
$L_{\mathrm{diff}}=44\,\mathrm{nm}$ (the value used in the example netlist provided with the assignment) for every transistor in every netlist; perimeters (PD/PS) are not set.
The model-file header mentions $0.95\,\mathrm{V}$; the assigned $1\,\mathrm{V}$ is used throughout.

\subsection{Repository and pipeline}
All results in this report are produced by the numbered scripts in \texttt{scripts/}, run in order (Appendix~\ref{app:reproduce}), using the library
\texttt{src/dvdml/}. The scripts read every assignment parameter from \texttt{config/project.yaml} and the model card from \texttt{22nm\_LP.pm}.
Every table in this report is generated from the result files by \texttt{scripts/11\_build\_report\_tables.py}; no number in a table is typed by hand.
The location of each deliverable is listed in Appendix~\ref{app:deliverables}.
